\subsection{Intención e Idioma}

\subsubsection{Regla: explicar el porqué en inglés}

\textbf{Regla:} El código debe ser autoexplicativo por su nombrado. Un comentario se admite únicamente cuando explica una razón de negocio, una restricción técnica o una decisión que el código no puede expresar por sí mismo. Todo comentario se escribe en inglés y nunca repite lo que hace la instrucción.

\textbf{Descripción:} Martin sostiene que la mayoría de los comentarios compensan código poco claro; el comentario valioso explica una intención que el código no puede comunicar por sí solo (Martin, 2008).

\textbf{Ejemplo correcto:} ``The official rules grant an extra turn after the maximum roll.'' explica una regla de negocio que no resulta evidente por la condición.

\textbf{Ejemplo incorrecto:} ``If the roll equals the maximum value, grant an extra turn.'' repite literalmente la condición y la llamada ejecutada.
